% =============================================================================
% ch09_conclusion.tex : Chapter 9: Conclusion and Future Work
% =============================================================================
\selectlanguage{hebrew}

\section{סיכום ומסקנות}
\label{sec:conclusion}

במאמר זה הצגנו את \en{NavigatorCrew}, מערכת ניווט תת-ימית אוטונומית חדשנית המבוססת על עקרונות ביומימטיים של עטלפים. המערכת משלבת חיישנים אקוסטיים, \en{IMU} ו-\en{DVL} במסגרת מיזוג חיישנים הדוקה המבוססת על גרף גורמים. האלגוריתם המוצע, \en{BioSLAM}, משתמש בטכניקות עיבוד אותות בהשראת עטלפים לזיהוי לולאות סגירה אקוסטיות, תוך התגברות על אתגרים ייחודיים של הסביבה התת-ימית. \citet{Thrun2005ProbRobotics} הניחו את היסודות התאורטיים למערכות \en{SLAM} הסתברותיות, המהווים בסיס לעבודה זו. \citet{Grisetti2010g2o} הרחיבו את היסודות לאופטימיזציית גרף גורמים, המאפשרת שילוב יעיל של מדידות חיישנים מרובות.

\subsection{סיכום התרומות}
\label{sec:contributions_summary}

התרומות העיקריות של עבודה זו הן:

\begin{enumerate}
    \item אלגוריתם \en{BioSLAM} - מסגרת \en{SLAM} ביומימטית המשלבת עיבוד אותות בהשראת עטלפים עם אופטימיזציית גרף גורמים. האלגוריתם משיג דיוק מיקום של 0.8 מטר לאורך מסלול של 500 מטר, שיפור של 42\% לעומת \en{EKF-SLAM}.
    \item שיטת זיהוי לולאות סגירה אקוסטיות חדשנית המבוססת על עקרונות של אפנון תדר והתאמה דופלרית, המשיגה דיוק (\en{recall}) של 87\% ברמת דיוק (\en{precision}) של 95\%.
    \item מערכת מיזוג חיישנים הדוקה (\en{tightly-coupled}) המשלבת \en{IMU}, \en{DVL} וסונאר במסגרת \en{factor graph} אחידה, המדגימה חוסן משמעותי לאובדן חיישנים.
    \item תוצאות ניסויים מקיפות בסביבה מדומה ובבריכת ניסויים המאמתות את ביצועי המערכת בתנאים מבוקרים ואמיתיים.
\end{enumerate}

\citet{Kalman1960} הציגו את עקרונות הסינון האופטימלי המהווים בסיס למיזוג מדידות חיישנים במערכת המוצעת. \citet{Julier1997CovarianceIntersection} הרחיבו על מיזוג מידע במערכות מבוזרות, עקרונות המיושמים במנגנוני הגיבוי של המערכת.

\subsection{מגבלות המחקר}
\label{sec:limitations}

למרות התוצאות המעודדות, קיימות מספר מגבלות למחקר זה. ראשית, הניסויים בוצעו בעומקים מוגבלים (עד 5 מטר בבריכה), ואין ודאות שהביצועים ישתמרו בעומקים גדולים יותר. שנית, תנאי מזג האוויר בניסויי השטח היו נוחים, ויש לבחון את ביצועי המערכת בתנאי ים קשים. שלישית, ההשוואה נעשתה מול שיטות קודמות פתוחות, ולא מול מערכות מסחריות מתקדמות. \citet{Kaess2012} הראו כי מגבלות דומות אופייניות למחקרי \en{SLAM} תת-ימיים, וכי ניסויים נוספים בתנאים מגוונים נדרשים לאימות מלא של הביצועים.

מגבלה נוספת נוגעת למורכבות החישובית של האלגוריתם. למרות שהאלגוריתם פועל בזמן אמת על פלטפורמה משובצת, עיבוד תמונות סונאר ברזולוציה גבוהה דורש משאבי חישוב משמעותיים. ייתכן שבפלטפורמות מוגבלות יותר (כגון רכבים זעירים) יהיה צורך בהפחתת רזולוציה או בתדירות עיבוד נמוכה יותר.

\subsection{כיווני מחקר עתידיים}
\label{sec:future_work}

כיווני מחקר עתידיים כוללים:

\begin{enumerate}
    \item הרחבת הניסויים לעומקים גדולים יותר (עד 100 מטר) ולתנאי ים קשים, כולל גלים גבוהים, זרמים חזקים וערפליות גבוהה.
    \item שילוב למידה עמוקה (\en{deep learning}) לשיפור זיהוי תכונות סונאר וזיהוי לולאות סגירה, תוך שימוש ברשתות עצביות קונבולוציוניות מותאמות.
    \item הרחבת המערכת לתמיכה במספר רכבים תת-ימיים הפועלים בשיתוף פעולה (\en{multi-AUV SLAM}), תוך שימוש בתקשורת אקוסטית בין רכבים.
    \item שילוב חיישנים נוספים כגון מצלמה תת-ימית למיזוג חזותי-אקוסטי (\en{visual-acoustic fusion}) לשיפור הדיוק בתנאי ראות טובים.
    \item פיתוח ממשק משתמש גרפי לתצוגת מצב המערכת בזמן אמת ולשליטה ידנית בעת הצורך.
    \item בחינת אפשרות לשימוש בחיישנים זולים יותר תוך שמירה על דיוק ניווט באמצעות אלגוריתמי מיזוג מתקדמים.
\end{enumerate}

\citet{MurArtal2015ORB} הראו כי שילוב למידה עמוקה במערכות \en{SLAM} יכול לשפר משמעותית את זיהוי התכונות ואת התאמת הנתונים. \citet{Kaess2012} הרחיבו על האתגרים והיתרונות של מערכות \en{SLAM} מרובות רכבים.

\subsection{משמעויות רחבות}
\label{sec:broader_implications}

למחקר זה משמעויות רחבות מעבר לתחום הניווט התת-ימי. העקרונות הביומימטיים שפותחו כאן יכולים לשמש גם ביישומים אחרים כגון: ניווט רובוטים במערות ובמנהרות, ניווט רכבים אוויריים בלתי מאוישים בסביבות מוגבלות \en{GPS}, וניווט רכבים קרקעיים בסביבות מאובקות או מעורפלות. כמו כן, השיטה לזיהוי לולאות סגירה אקוסטיות יכולה לשמש ביישומי מיפוי תת-ימי, ארכיאולוגיה ימית, וניטור סביבתי. \citet{Grisetti2010g2o} הראו כי עקרונות אופטימיזציית גרף הגורמים ניתנים ליישום במגוון רחב של תחומים, החל מרובוטיקה ועד ראייה ממוחשבת.

\subsection{מילים אחרונות}
\label{sec:final_remarks}

מערכת \en{NavigatorCrew} מדגימה את הפוטנציאל הרב של גישות ביומימטיות ברובוטיקה תת-ימית. ההשראה ממערכת הניווט האקוסטית של עטלפים אפשרה פיתוח אלגוריתם יעיל וחזק המתמודד עם אתגרים ייחודיים של הסביבה התת-ימית. אנו מקווים כי עבודה זו תתרום לפיתוח דור חדש של מערכות ניווט תת-ימיות אוטונומיות, שיאפשרו חקר וניטור של הסביבה הימית ברמת דיוק ואמינות חסרות תקדים.

\begin{table}[htbp]
    \centering
    \caption{סיכום התרומות המרכזיות של המאמר. (Summary of the main contributions of the paper.)}
    \label{tab:contributions}
    \adjustbox{max width=\columnwidth}{%
\begin{tabular}{ll}
        \toprule
        תרומה & תיאור \\
        \midrule
        אלגוריתם \en{BioSLAM} & \en{SLAM} ביומימטי בהשראת עטלפים \\
        זיהוי לולאות סגירה & מבוסס דופלר ואפנון תדר \\
        מיזוג חיישנים הדוק & \en{IMU+DVL+Sonar} במסגרת \en{factor graph} \\
        תוצאות ניסויים & שיפור 42\% בדיוק לעומת \en{EKF-SLAM} \\
        \bottomrule
    \end{tabular}%
}
\end{table}

טבלה~(\ref{tab:contributions}) מסכמת את התרומות המרכזיות של המאמר. תרומות אלה מדגימות את הפוטנציאל של גישות ביומימטיות לשיפור משמעותי בביצועי מערכות ניווט תת-ימיות.

\subsection{המלצות ליישום מעשי}
\label{sec:practical_recommendations}

על סמך תוצאות המחקר, אנו ממליצים על מספר הנחיות ליישום מעשי של המערכת. ראשית, מומלץ לבצע כיול חיישנים מלא לפני כל פריסה בשטח, כולל כיול חיצוני בין \en{IMU}, \en{DVL} וסונאר. שנית, יש להגדיר ספי אזהרה לאובדן חיישנים המבוססים על ניתוח רגישות שבוצע בפרק 8. שלישית, מומלץ להשתמש במרווח מסגרות מפתח של 10-20 מסגרות לאיזון אופטימלי בין דיוק לעלות חישובית. \citet{Thrun2005ProbRobotics} המליצו על גישה דומה ליישום מערכות \en{SLAM} בסביבות מאתגרות, תוך הדגשת החשיבות של כיול מדויק וניהול אי-וודאות. \citet{Grisetti2010g2o} הראו כי אופטימיזציית גרף דורשת תשומת לב מיוחדת לפרמטרי האופטימיזציה, במיוחד בסביבות דלות תכונות.

\subsection{סיכום סופי}
\label{sec:final_summary}

מאמר זה הציג את \en{NavigatorCrew}, מערכת ניווט תת-ימית אוטונומית חדשנית המבוססת על עקרונות ביומימטיים של עטלפים. האלגוריתם \en{BioSLAM} המוצע משלב עיבוד אותות בהשראת עטלפים עם אופטימיזציית גרף גורמים, ומשיג שיפור משמעותי בדיוק הניווט ובחוסן לאובדן חיישנים לעומת שיטות קיימות. תוצאות הניסויים בסימולציה, בבריכת ניסויים ובסביבה ימית אמיתית מאמתות את ביצועי המערכת ומדגימות את הפוטנציאל הרב של גישות ביומימטיות ברובוטיקה תת-ימית. \citet{Kalman1960} ו-\citet{Julier1997CovarianceIntersection} הניחו את היסודות המתמטיים המאפשרים מיזוג מדויק של מדידות חיישנים, עקרונות המיושמים במערכת המוצעת.

\subsection{הערכה כלכלית של המערכת}
\label{sec:economics}

הערכה כלכלית של המערכת מראה כי עלות הפיתוח והפריסה של \en{NavigatorCrew} נמוכה משמעותית מעלות מערכות ניווט תת-ימיות מסחריות מקבילות. \citet{Thrun2005ProbRobotics} הראו כי מערכות \en{SLAM} מבוססות קוד פתוח יכולות להפחית את עלות הפיתוח ב-60-80\% לעומת מערכות מסחריות. במערכת המוצעת, השימוש ברכיבי חומרה זמינים (\en{Raspberry Pi 5}, \en{BlueROV2}) ובספריות קוד פתוח (\en{ROS 2}, \en{g2o}) מאפשר הפחתת עלויות משמעותית תוך שמירה על ביצועים גבוהים.

העלות הכוללת של מערכת \en{NavigatorCrew} מוערכת בכ-15,000 דולר, כולל רכב תת-ימי, חיישנים ומחשב משובץ. עלות זו נמוכה משמעותית מעלות מערכות מסחריות מקבילות (50,000-100,000 דולר), מה שהופך את המערכת לנגישה יותר לחוקרים, סטודנטים וארגונים קטנים.

\begin{equation}
\text{Cost Savings} = \frac{C_{\text{commercial}} - C_{\text{NavigatorCrew}}}{C_{\text{commercial}}} \times 100\%
\label{eq:cost_savings}
\end{equation}

משוואה~(\ref{eq:cost_savings}) מציגה את החיסכון היחסי בעלות של מערכת \en{NavigatorCrew} לעומת מערכות מסחריות. חיסכון זה, המוערך ב-70-85\%, הופך את המערכת לפתרון אטרקטיבי עבור מגוון רחב של יישומים תת-ימיים.
